\documentclass[11pt]{article}

% some definitions for the title page
\newcommand{\reporttitle}{What is Radiotherapy?}
\newcommand{\reportdescription}{}

% load some definitions and default packages
%%%%%%%%%%%%%%%%%%%%%%%%%%%%%%%%%%%%%%%%%
% University Assignment Title Page 
% LaTeX Template
% Version 1.0 (27/12/12)
%
% This template has been downloaded from:
% http://www.LaTeXTemplates.com
%
% Original author:
% WikiBooks (http://en.wikibooks.org/wiki/LaTeX/Title_Creation)
%
% License:
% CC BY-NC-SA 3.0 (http://creativecommons.org/licenses/by-nc-sa/3.0/)
% 
%
%%%%%%%%%%%%%%%%%%%%%%%%%%%%%%%%%%%%%%%%%
%----------------------------------------------------------------------------------------
%	PACKAGES AND OTHER DOCUMENT CONFIGURATIONS
%----------------------------------------------------------------------------------------
\usepackage[a4paper,hmargin=2.0cm,vmargin=1.0cm,includeheadfoot]{geometry}
\usepackage{textpos}

\usepackage[square,numbers]{natbib} % for bibliography
\usepackage[nottoc]{tocbibind} % Includes "References" in the table of contents
\bibliographystyle{unsrtnat}

\usepackage{tabularx,longtable,multirow,subfigure,caption}%hangcaption
\usepackage{fancyhdr} % page layout
\usepackage{url} % URLs
\usepackage[english]{babel}
\usepackage{amsmath}
\usepackage{graphicx}
\usepackage{scalerel}
\usepackage{dsfont}
\usepackage{epstopdf} % automatically replace .eps with .pdf in graphics
\usepackage{backref} % needed for citations
\usepackage{array}
\usepackage{latexsym}

\usepackage[pdftex,pagebackref,hypertexnames=false,colorlinks]{hyperref} % provide links in pdf
\usepackage{booktabs}
\usepackage{wrapfig}
\usepackage{caption}  % Required for \captionof
\usepackage{float} % for H option in figures
\usepackage{amssymb}
\usepackage{amsmath}
\usepackage{csquotes}
% \usepackage{subcaption} % causes a compilation error after changing back to natbib referencing... 

\hypersetup{pdftitle={},
  pdfsubject={}, 
  pdfauthor={},
  pdfkeywords={}, 
  pdfstartview=FitH,
  pdfpagemode={UseOutlines},% None, FullScreen, UseOutlines
  bookmarksnumbered=true, bookmarksopen=true, colorlinks,
    citecolor=black,%
    filecolor=black,%
    linkcolor=black,%
    urlcolor=black}

\usepackage[all]{hypcap}

%%%%%%%%%%%%%%%%%%%%%%%%%%%%%%%%%%%%%%%%%%%%%%%%%%%%%%%%%%%%%%%%%%%%%%%%%%%%%%%%
% LISTINGS ammendments
%%%%%%%%%%%%%%%%%%%%%%%%%%%%%%%%%%%%%%%%%%%%%%%%%%%%%%%%%%%%%%%%%%%%%%%%%%%%%%%%
\usepackage{listings}

\lstset
{ %Formatting for code in appendix
    language=Matlab,
    basicstyle=\footnotesize,
    % numbers=left,
    stepnumber=1,
    showstringspaces=false,
    tabsize=1,
    breaklines=true,
    breakatwhitespace=false,
    frame=single,
    columns=fullflexible,
    postbreak=\mbox{\textcolor{red}{$\hookrightarrow$}\space},
}

%\usepackage{color}
%\usepackage[tight,ugly]{units}
%\usepackage{float}
%\usepackage{tcolorbox}
%\usepackage[colorinlistoftodos]{todonotes}
% \usepackage{ntheorem}
% \theoremstyle{break}
% \newtheorem{lemma}{Lemma}
% \newtheorem{theorem}{Theorem}
% \newtheorem{remark}{Remark}
% \newtheorem{definition}{Definition}
% \newtheorem{proof}{Proof}


%%% Default fonts
\renewcommand*{\rmdefault}{bch}
\renewcommand*{\ttdefault}{cmtt}


%%% Default settings (page layout)
\setlength{\parindent}{0em}  % indentation of paragraph
\setlength{\parskip}{.3em}

% \setlength{\parindent}{0em}  % indentation of paragraph

\setlength{\headheight}{14.5pt}
\pagestyle{fancy}
\renewcommand{\chaptermark}[1]{\markboth{\chaptername\ \thechapter.\ #1}{}}
%\fancyhead[RO]{\sffamily \textbf{\thepage}} %Page no.in the right on even pages
%\fancyhead[LE]{\sffamily \textbf{\thepage}} %Page no. in the left on odd pages

\fancyfoot[ER,OL]{\thepage}%Page no. in the left on
%odd pages and on right on even pages
\fancyfoot[OC,EC]{\sffamily }
\renewcommand{\headrulewidth}{0.1pt}
\renewcommand{\footrulewidth}{0.1pt}
\captionsetup{margin=10pt,font=small,labelfont=bf}


%--- chapter heading

\def\@makechapterhead#1{%
  \vspace*{10\p@}%
  {\parindent \z@ \raggedright \sffamily
    \interlinepenalty\@M
    \Huge\bfseries \thechapter \space\space #1\par\nobreak
    \vskip 30\p@
  }}

%--- chapter heading

\def\@makechapterhead#1{%
  \vspace*{10\p@}%
  {\parindent \z@ \raggedright \sffamily
    %{\Large \MakeUppercase{\@chapapp} \space \thechapter}
    %\\
    %\hrulefill
    %\par\nobreak
    %\vskip 10\p@
    \interlinepenalty\@M
    \Huge\bfseries \thechapter \space\space #1\par\nobreak
    \vskip 30\p@
  }}

%---chapter heading for \chapter*  
\def\@makeschapterhead#1{%
  \vspace*{10\p@}%
  {\parindent \z@ \raggedright
    \sffamily
    \interlinepenalty\@M
    \Huge \bfseries  #1\par\nobreak
    \vskip 30\p@
  }}
\allowdisplaybreaks
% Here, you can define your own macros. Some examples are given below.

\newcommand{\R}[0]{\mathds{R}} % real numbers
\newcommand{\Z}[0]{\mathds{Z}} % integers
\newcommand{\N}[0]{\mathds{N}} % natural numbers
\newcommand{\C}[0]{\mathds{C}} % complex numbers
% \renewcommand{\vec}[1]{{\boldsymbol{{#1}}}} % vector
\newcommand{\mat}[1]{{\boldsymbol{{#1}}}} % matrix

\usepackage{pifont,mdframed}
\newenvironment{warning}
  {\par\begin{mdframed}[linewidth=1pt,linecolor=black]%
    \begin{list}{}{\leftmargin=1cm
                  \labelwidth=\leftmargin}\item[\Large\ding{43}]}
  {\end{list}\end{mdframed}\par}

\definecolor{lightgray}{gray}{0.9}

\begin{document}

% Include the title page
% Last modification: 2015-08-17 (Marc Deisenroth)
\begin{titlepage}

    \newcommand{\HRule}{\rule{\linewidth}{0.5mm}} % Defines a new command for the horizontal lines, change thickness here
    
    %----------------------------------------------------------------------------------------
    %	LOGO SECTION
    %----------------------------------------------------------------------------------------
    
    \includegraphics[width = 4cm]{../figures/imperial.pdf}\\[0.5cm] 
    
    \center % Center everything on the page
     
    %----------------------------------------------------------------------------------------
    %	HEADING SECTIONS
    %----------------------------------------------------------------------------------------
    
    \textsc{\LARGE \reporttype}\\[1.5cm] 
    \textsc{\Large Department of Computing}\\[0.5cm] 
    \textsc{\large Imperial College of Science, Technology and Medicine}\\[0.5cm] 
    
    %----------------------------------------------------------------------------------------
    %	TITLE SECTION
    %----------------------------------------------------------------------------------------
    
    \HRule \\[0.4cm]
    { \huge \bfseries \reporttitle}\\ % Title of your document
    \HRule \\[1.5cm]
     
    %----------------------------------------------------------------------------------------
    %	AUTHOR SECTION
    %----------------------------------------------------------------------------------------
    
    \begin{minipage}{0.4\textwidth}
    \begin{flushleft} \large
    \emph{Author:}\\
    \reportauthor % Your name
    \end{flushleft}
    \end{minipage}
    ~
    \begin{minipage}{0.4\textwidth}
    \begin{flushright} \large
    \emph{Supervisor:} \\
    \supervisor % Supervisor's Name
    \end{flushright}
    \end{minipage}\\[4cm]
    
    
    
    
    %----------------------------------------------------------------------------------------
    
    
    %----------------------------------------------------------------------------------------
    %	DATE SECTION
    %----------------------------------------------------------------------------------------
    
    {\large \today} % Date, change the \today to a set date if you want to be precise
    
    
    \vfill % Fill the rest of the page with whitespace
    Submitted in partial fulfillment of the requirements for the \degreetype~of Imperial College London
    
    \end{titlepage}
    

\tableofcontents

\clearpage

\section{What is Radiotherapy?}

Clinicians use a system that uses high-resolution X-rays to produce contrasting images of cancerous tumours and surrounding soft tissue. Physicians then can target the cancerous tumour more precisely and decreasing radiation expsoure of healthy tissues~\cite{radiotherapy-basic-concepts}.

Radiation therapy (radiotherapy) remains an improtant component of cancer treatment. In 2012 approximately 50\% of all cancer patients received radiation therapy, with an additional 40\% involving curative treatment~\cite{radiotherapy-advances}.

\subsection{Cell Dealth}

Radiation-induced cell death is typically classified as interphase or proliferative death. Importantly, there is a consensus that cellular effects, including cell death, depend not only on radiation dose but cell type and a place in cell cycle~\cite{cell-death}.

\subsubsection{Interphase death}

Cells cease to divide after radiation exposure and begin to die within hours. ``This is attributed to the damage of intercellular molecules and the activation of nuclease and proteolytic enzymes etc., after high doses of radiation, which leads to degradation [...]. Factors such as disruption of membrane structure and disorder of cell energy metabolism after irradiation are also important contributors to interphase death''~\cite{cell-death}.

\subsubsection{Proliferative death}

Most cells undergo proliferative death after radiation. This results from ``mitotic catastrophe caused by accumulation of chromosomal aberrations and erroneous repair after radiation induces a DNA double-strand break [...] cells lose their ability to proliferate and begin to die.''~\cite{cell-death}.

\subsection{Using Radiotherapy for Cancer}

Physicians use high-energy radiation to damage genetic material (deoxyribonucleic acid, DNA) of cells and thus block their ability to divide and proliferate further~\cite{radiotherapy-advances}.

\subsection{The need for Radiotherapy}

% https://radiotherapy.org.uk/wp-content/uploads/2022/11/APPGRadiotherapy_Manifesto-Update_2022.pdf

`In England, waiting times are getting worse each month – over 1/3 of cancer patients (60,000) wait beyond the 62-day target, and 10,000 patients wait over 104 days'.
 
\section{Cervical Cancer}

In 2017, Cervical Cancer was the fourth most common malignancy diagnosed in women world wide~\cite{cervical-cancer-epidemic}.

\begin{itemize}
    \item There are 530,000 new cases annaully, with 270,000 deaths~\cite{cervical-cancer-epidemic}.
    \item approximately 85\% of cases occurr in underdeveloped or developing countries~\cite{cervical-cancer-epidemic}
    \item The median age for diagnosis is 47~\cite{cervical-cancer-epidemic} which is by no means an end-of-life diagnosis, and many victims will still have a lot of time remaining after proper treatment.
    \item `The standard management of individuals with locally advanced cervical cancer includes external beam radiotherapy'~\cite{cervical-cancer-epidemic}.
    \item `Scale-up of radiotherapy capacity in 2015-35 from current levels could lead to saving of 26.9 million life-years in low-income and middle-income countries over the lifetime of the patients who received treatment'~\cite{expanding-global-access-to-radiotherapy}
\end{itemize}

\printbibliography

\end{document}